\documentclass[11pt]{article}
\usepackage[T1]{fontenc}
\usepackage{lmodern,microtype}
\usepackage[letterpaper,margin=0.86in]{geometry}
\usepackage{amsmath,amssymb,amsthm,mathtools,booktabs,longtable,array}
\usepackage{xcolor,enumitem,fancyhdr}
\usepackage[breaklinks,colorlinks=true,linkcolor=blue!45!black,citecolor=blue!45!black,urlcolor=blue!45!black]{hyperref}
\definecolor{navy}{RGB}{24,43,68}
\definecolor{teal}{RGB}{18,105,116}
\setlength{\parindent}{0pt}
\setlength{\parskip}{5pt}
\setlength{\headheight}{14pt}
\pagestyle{fancy}\fancyhf{}
\fancyhead[L]{\small\color{navy}Thirteen integer variables for Con(ZFC)}
\fancyhead[R]{\small 23 October 2026}
\fancyfoot[C]{\thepage}
\setlist{nosep,leftmargin=1.5em}
\allowdisplaybreaks[2]
\newtheorem{theorem}{Theorem}[section]
\newtheorem{lemma}[theorem]{Lemma}
\newtheorem{proposition}[theorem]{Proposition}
\newtheorem{corollary}[theorem]{Corollary}
\theoremstyle{remark}\newtheorem{remark}[theorem]{Remark}
\newcommand{\Z}{\mathbb Z}\newcommand{\Q}{\mathbb Q}\newcommand{\N}{\mathbb N}
\newcommand{\ZFC}{\mathsf{ZFC}}\newcommand{\Con}{\operatorname{Con}}
\newcommand{\Pos}{\operatorname{Pos}}\newcommand{\ord}{\operatorname{ord}}
\newcommand{\file}[1]{\path{#1}}
\newcommand{\pin}{\texttt{cfbf118b4ad9b2c651fe96b571d9415261766328}}
\hypersetup{pdftitle={Thirteen integer variables for Diophantine ZFC consistency},pdfauthor={Prepared for Vladimir Reshetnikov; AI-assisted mathematical research},pdfsubject={Fixed-prime radical compression, Pell equations and syntactic consistency}}

\begin{document}
\thispagestyle{empty}
{\sffamily\bfseries\color{teal} RESEARCH CONTRIBUTION \hfill 23 OCTOBER 2026\par}
\vspace{0.85cm}
{\Huge\bfseries\color{navy} Thirteen integer variables\\[4pt] for Diophantine\\[4pt] ZFC consistency\par}
\vspace{0.5cm}
{\Large Fixed-prime radical compression and a compatible positivity test\par}
\vspace{0.7cm}
Prepared for Vladimir Reshetnikov and the ProveIt project.\\
AI-assisted research with ordinary mathematical proofs, independently reviewed\\
reductions, literal polynomial circuits, fixed coefficients and exact checks.
\vspace{0.65cm}
\begin{abstract}
We give a fixed ordinary integer-coefficient polynomial in \textbf{thirteen
unrestricted integer unknowns}, of exact total degree \textbf{1392}, with a
\textbf{475-operation} arithmetic circuit. Composing the reductions proved
here with the explicitly retained ProveIt native/compiler calibration proves
in ordinary mathematics formalizable in ZFC that its integer solvability is
equivalent to a finite syntactic ZFC refutation. Negating gives the requested
equivalence between its lack of integer zeros and $\Con(\ZFC)$.

The central mechanism uses distinct valuations at the prime $3$ to combine
several square conditions with fixed small weights. This replaces the
size-dependent separation used by the inspected predecessor constructions.
A normalized strong Pell witness makes a fourth root available for
compression. A classical three-pentagonal representation gives a positivity
radicand that is always $1$ modulo $3$, allowing a fifth root to join the same
norm. The complete norm has degree $32$ in its carrier; every radical and
division disappears from the supplied polynomial circuit.

A companion uses \textbf{fourteen unrestricted integers}, exact degree
\textbf{652}, and \textbf{230 operations}. We retain uncentered and partially
centered control equations so that every degree-saving coordinate change is
explicit. The six inherited numerical coefficient payloads total 615,883
uncompressed bytes and are reconstructed byte for byte.

This is a substantial new degree/variable tradeoff within the reviewed
repository, not a proof of global simplicity or a new fewest-variable
theorem. Ten-integer equations already exist at much higher degree. The new
arithmetic proofs and finite verification are distinct from an end-to-end
Lean/Rocq certificate or an exported derivation in ProveIt's encoded ZFC
calculus; neither is claimed here.
\end{abstract}
\vfill
\small Comparison snapshot: \pin.\\
All numbered packages 31--75 and the incoming research archives at that
snapshot were included in the comparison.
\newpage
\begingroup\small\setlength{\parskip}{0pt}\tableofcontents\endgroup
\newpage

\section{Result and comparison}

\subsection{The logical statement}
Fix the recursive syntactic ZFC refutation predicate used by the retained
ProveIt calibration. Its consistency sentence says that no finite refutation
exists. The result is a fixed polynomial $P_{13}\in\Z[x_1,\ldots,x_{13}]$
with the following theorem:
\begin{equation}\label{eq:target}
 \ZFC\vdash
 \left(\neg\exists\mathbf x\in\Z^{13}\;P_{13}(\mathbf x)=0\right)
 \longleftrightarrow\Con(\ZFC).
\end{equation}
The proof first establishes the existence equivalence with a finite
refutation and then negates it. It assumes neither consistency nor
inconsistency. Integer quantification is the usual definable arithmetic
quantification inside ZFC, including negative integers and zero.

The exact new arithmetic claim is independent of any assertion about ZFC:
the new polynomial has a zero if and only if the retained calibrated source
has one. The last source-to-refutation step uses named ordinary theorems
from the pinned project. Their source texts, the exact parent polynomial,
the literal machine and the coefficient recipe are retained in the archive.
Section~\ref{sec:scope} distinguishes this mathematical composition from the
remaining proof-assistant and object-language work.

\subsection{Measured equations}
Every operation in this table is one live binary addition, subtraction or
multiplication. Squaring and multiplication by a fixed numeral count;
numeral leaves and register reuse are free. Thus these are arithmetic-circuit
counts, not bit-complexity estimates.
\begin{center}\small
\begin{tabular}{@{}lrrr@{}}\toprule
Unrestricted-integer equation & Unknowns & Exact degree & Operations\\\midrule
New main equation, auxiliary root centered & 13 & \textbf{1392} & \textbf{475}\\
New uncentered control & 13 & 1408 & 473\\
New companion, three roots centered & 14 & \textbf{652} & \textbf{230}\\
New integer-center control & 14 & 664 & 223\\
New auxiliary-center control & 14 & 700 & 217\\
New uncentered control & 14 & 748 & 215\\\bottomrule
\end{tabular}
\end{center}
An additional companion retains the ordinary strong root: it uses
\textbf{15 unrestricted integers, degree 214, and 139 operations}.
All seven equations have the same integer-zero existence status. Their
different numbers of private positivity coordinates are stated explicitly;
the named registers in their circuits are not additional unknowns.

\subsection{What was already available}
The following are selected previously delivered points from the actual
comparison snapshot, including the incoming archives. They must not be
replaced by the older comparison table in package 75.
\begin{center}\small
\begin{tabular}{@{}llrr@{}}\toprule
Source & Domain & Unknowns & Degree\\\midrule
Structural Round11 & $\Z$ & 37 & 4\\
Structural Round11 & $\Z$ & 30 & 6\\
Joint Guards & $\Z$ & 25 & 8\\
Polynomial Shears & $\Z$ & 23 & 10\\
Centered Pell Norms & $\Z$ & 22 & 14\\
Polynomial Shears & $\Z$ & 21 & 16\\
Pell Congruences & $\Z$ & 20 & 20\\
Pell Congruences & $\Z$ & 19 & 28\\
Pell Congruences & $\Z$ & 18 & 32\\
Pell Congruences, immediate parent & $\Z$ & 17 & 88\\
Package 67, named thirteen-integer form & $\Z$ & 13 & 10,461,744\\
Polynomial Shears, few-variable alternative & $\Z$ & 10 & 1,081,125\\\midrule
Radical Residue & $\Z_{>0}$ & 10 & 1152\\
Polynomial Shears & $\Z_{>0}$ & 11 & 211\\
Saturation and Projection, lower-degree nine & $\Z_{>0}$ & 9 & 1,103,950\\\bottomrule
\end{tabular}
\end{center}
The new thirteen-integer form lowers the degree of the named thirteen-integer
predecessor by a factor exceeding 7500. The existing ten-integer forms remain
better in variable count; the existing seventeen-integer form remains better
in degree. Domain restrictions are material: the ten-positive degree-1152
equation is not a ten-unrestricted-integer equation.

The new mechanism was not found in the inspected repository or incoming
texts. We make no worldwide priority claim. The pentagonal theorem used
below is classical and is expressly attributed. The inherited native
encoding and K215 calibration are retained; a completely different ZFC
verifier has not been compiled in this contribution.

\section{The complete thirteen-variable equation}\label{sec:equation}

\subsection{Variables and fixed coefficients}
The independent variables are exactly
\begin{equation}\label{eq:vars}
 j,C,z,h,r,v,w,T,c,s,u_1,u_2,\xi\in\Z.
\end{equation}
All other symbols below are finite polynomial abbreviations in these
thirteen integers. The six fixed coefficients are
\[
 (\beta,K_0,M_C,M_F,A_0,B_0),\qquad
 \beta=2^{1064921}-1,\quad A_0=2^{1113},\quad B_0=2^{666104}.
\]
The middle three integers belong to the complete source34
\emph{mixed-parity-unrounded-1064921} tuple. Their exact byte values and
finite constructive definition are included in \file{coefficients/}.
$M_F$ is already shifted by $\beta$. These are fixed numerals, not
parameters quantified along with the unknowns. No coefficient depends on
waiting for a halting computation.
The literal JSON names $u_1,u_2$ as \texttt{p1,p2}, respectively, and
names $\xi$ as \texttt{xi}; these are only typographical renamings.

\subsection{Native and Pell abbreviations}
Define, in this order,
\begin{align}
 J&=6j+1,& q&=\beta J+1,& W&=4(3w+1),\notag\\
 F&=(K_0+W)C-z(q-1),& X&=qW,&Y&=(2v+1)q^3,\notag\\
 R&=\bigl(q(q-F)-C\bigr)(q^2-1)+(M_C+qM_F)J,\label{eq:prefix}\\
 E&=XY,&a&=E+Y,&H&=4a+3,\notag\\
 \Delta&=a^2+H=(a+2)^2-1,&k&=R+1+hE,\notag\\
 \eta&=c-kY,&\zeta&=k-\eta,&\ell&=EkY.\notag
\end{align}
The normalized strong and auxiliary expressions are
\begin{align}
 I_0&=(2r+1)c^2,& D_f&=1+\Delta I_0^2,\notag\\
 Q&=\Delta^2I_0^2,&\delta&=Q-1,\notag\\
 V&=D_f(cT-R)-c,&D_0&=X+ac,\label{eq:pell}\\
 A_\tau&=\ell(\ell+k)+1,&D_m&=1+\Delta c^2.\notag
\end{align}
The two sign quantities and their joint positivity test are
\begin{align}
 G&=A_0FB_0C-(q-A_0F-B_0C)^2,\notag\\
 Z&=\eta\zeta,\notag\\
 N&=GZ-(s^2-G-Z+1)^2,\label{eq:guard}\\
 S&=24N-21-(6u_1-1)^2-(6u_2-1)^2.\notag
\end{align}
Finally set
\begin{align}
 p&=D_0+H(\xi+27V),\notag\\
 A&=9H^2A_\tau,& B&=81H^2D_f,& D&=D_m,\label{eq:normports}\\
 L&=6561H^2S,& C_*&=729H^2(QV^2-1).\notag
\end{align}
The letter $L$ here is only a radicand abbreviation; it is unrelated to the
native block length in the coefficient generator.

\subsection{A finite division-free norm definition}
For $m\ge0$ define the universal polynomial
\begin{equation}\label{eq:universalnorm}
 \mathcal F_m(x;d_1,\ldots,d_m)=
 \prod_{\varepsilon\in\{-1,1\}^m}
 \left(x-\sum_{i=1}^m\varepsilon_i\sqrt{d_i}\right),
 \qquad\mathcal F_0(x)=x.
\end{equation}
The product is invariant under changing each radical's sign, so every
radical occurs only to an even power. Therefore
$\mathcal F_m\in\Z[x,d_1,\ldots,d_m]$. Its degree in $x$ is $2^m$.
This is a definition of an ordinary polynomial, not an instruction to use
radicals in the final equation.

For a completely division-free finite prescription, expand
\[
 \mathcal F_4(p+t;A,B,D,L)=\sum_{j=0}^{16}b_jt^j.
\]
Every $b_j$ is an integer polynomial in $p,A,B,D,L$. Define
\begin{align}
 U&=\sum_{j=0}^{8} b_{2j} C_*^j\delta^{8-j},\notag\\
 V_*&=\sum_{j=0}^{7} b_{2j+1} C_*^j\delta^{7-j}.\label{eq:clear}
\end{align}
The promised equation is
\begin{equation}\label{eq:P13}
 \boxed{\quad P_{13}=U^2-\delta C_*V_*^2=0.\quad}
\end{equation}
Only finite sums, integer constants, additions, subtractions and fixed
nonnegative powers occur. No new unknowns are hidden in $U,V_*$ or in the
$b_j$. The supplied 475-operation graph expands this prescription into
binary arithmetic instructions.

For proof, where $\delta\ne0$, the defining identity is
\begin{equation}\label{eq:normidentity}
 P_{13}=\delta^{16}\mathcal F_5
       (p;A,B,D,L,C_*/\delta).
\end{equation}
The uncentered control replaces $p$ by $D_0+H\xi$.
This changes only the coordinate $\xi$: the map
$\xi\mapsto\xi+27V$ is an integral polynomial bijection, since $V$ does
not contain $\xi$.

\begin{theorem}\label{thm:main}
For the six fixed numerals specified above, $P_{13}$ has an integer zero
if and only if the retained source34 native construction has an accepting
certificate. By its retained ordinary calibration, this is equivalent to a
finite syntactic ZFC refutation. The exact degree of $P_{13}$ is 1392 and
the supplied circuit has 475 live binary operations.
\end{theorem}

\section{Fixed-prime separation of square roots}\label{sec:valuation}
The following elementary lemma is the central new tool in this construction.
It is deliberately stated for signed radicands. Before the guard is
recovered, two radicands could be negative, so an argument that merely counts
imaginary summands would be insufficient.

\begin{lemma}[Valuation-separated rational roots]\label{lem:separation}
Let $d_1,\ldots,d_m\in\Q^\times$. Suppose their valuations at a prime
$p_0$ are distinct even integers. If, for some choices of square roots in
$\mathbb C$, their signed sum is rational, then every $d_i$ is a square
in $\Q$.
\end{lemma}
\begin{proof}
Group the radicands by their classes in $\Q^\times/(\Q^\times)^2$.
For a class represented by $d$, write $\sqrt{d_i}=a_i\sqrt d$ with
$a_i\in\Q^\times$, absorbing signs into $a_i$. Since
\[
 2v_{p_0}(a_i)=v_{p_0}(d_i)-v_{p_0}(d),
\]
the $a_i$ in a class have pairwise distinct valuations. Their signed sum
cannot vanish: after dividing by its unique term of smallest valuation,
one obtains a rational unit plus terms divisible by $p_0$.

Distinct square classes yield linearly independent square roots over
$\Q$. One way to see this is to pass to the finite multiquadratic extension
generated by the roots and average a relation against its distinct
quadratic characters. Thus a rational total sum requires the coefficient
of each nontrivial square class to vanish. We just proved that no nonempty
class can have zero coefficient. Only the trivial square class remains.
The argument includes negative square classes and uses no order on the
radicands.
\end{proof}

\begin{corollary}\label{cor:norm}
If the $d_i$ satisfy the lemma, then
$\mathcal F_m(x;d_1,\ldots,d_m)=0$ with rational $x$ implies that all
$d_i$ are rational squares.
\end{corollary}
\begin{proof}
Over a splitting field the norm is the displayed finite product. A zero
factor gives the rational signed sum needed by the lemma.
\end{proof}

The weights need not make any root large in the usual absolute value.
They need only give distinct prime valuations inside every possible square
class. This is what permits small fixed multipliers $3,9,27,81$ even when
the Pell roots vary over arbitrarily large sizes.

\subsection{Why the five radicands are separated}
The literal definitions give, for every integer assignment,
\[
 q\equiv2\pmod3,\quad W\equiv1\pmod3,\quad
 X\equiv2\pmod3,\quad a\equiv\Delta\equiv H\equiv0\pmod3.
\]
Put $b=XY^2$. Modulo $3$, $b$ is either $0$ or $2$, so
\[
 A_\tau=1+b(b+1)k^2\equiv1\pmod3.
\]
Also $D_f,D_m,S\equiv1\pmod3$, while $Q\equiv0\pmod3$ and
\[
 B_y:=\frac{QV^2-1}{Q-1}\equiv1\pmod3.
\]
In particular $\delta=Q-1\ne0$ \emph{before} any positivity condition
has been recovered. This justifies denominator clearing without a
circular appeal to the guard.

Let $e=v_3(H)\ge1$. As $H=4a+3$ cannot vanish for integral $a$, the
five weighted squared roots in~\eqref{eq:normidentity} have valuations
\[
 2e+2,\qquad2e+4,\qquad0,\qquad2e+8,\qquad2e+6.
\]
They are pairwise distinct even integers. Hence a zero of $P_{13}$ makes
\[
 A_\tau,\quad D_f,\quad D_m,\quad S,\quad B_y
\]
all rational squares. The first four are integers and therefore have
integer square roots. The auxiliary rational square requires the descent
in Section~\ref{sec:descent}.

\section{A positivity test compatible with the same prime}\label{sec:positivity}

\begin{lemma}[Joint positivity]\label{lem:joint}
For integers $x,y$,
\[
 x>0\text{ and }y>0
 \quad\Longleftrightarrow\quad
 \exists s\in\Z:\;xy-(s^2-x-y+1)^2>0.
\]
\end{lemma}
\begin{proof}
Positivity of the expression gives $xy>0$. If both arguments were
negative, then $s^2-x-y+1\ge |x|+|y|+1>\sqrt{|xy|}$, a contradiction.
For the converse put $n=x+y-1\ge1$. Then $n\le xy$. Choose a square
nearest to $n$. Its squared error is less than $n$: between $a^2$ and
$(a+1)^2$, with integral $n$, the nearest error is at most $a$ and
$n>a^2$, unless $n$ is itself a square, when the error is zero. Thus the
required expression is positive.
\end{proof}

\begin{lemma}[A $3$-adic-unit positivity radicand]\label{lem:pentagonal}
For any integer $N$,
\[
 N>0\quad\Longleftrightarrow\quad
 \exists u_1,u_2\in\Z:\;
 24N-21-(6u_1-1)^2-(6u_2-1)^2
 \text{ is a square in }\Q.
\]
The displayed radicand is $1$ modulo $3$ for every assignment.
\end{lemma}
\begin{proof}
An integer that is a rational square is an integer square. If the
radicand equals $t^2$, then
$24N-21=(6u_1-1)^2+(6u_2-1)^2+t^2\ge2$, so $N\ge1$.
The residue assertion follows immediately.

For completeness put $M=24N-21=24(N-1)+3$. Remove the largest power
$9^e$ dividing $M$, obtaining $M_0$. We have $M_0>0$, $M_0\equiv3$
modulo $8$, and $9\nmid M_0$. Legendre's three-square theorem supplies
$M_0=x^2+y^2+z^2$. All three coordinates are odd and they are not all
divisible by $3$.

Use the integer matrix
\[
 L_0=\begin{pmatrix}1&2&2\\2&1&-2\\-2&2&-1\end{pmatrix},
 \qquad L_0^{\mathsf T}L_0=9I_3.
\]
For a vector nonzero modulo $3$, independent sign choices make
$x-y-z\ne0$ modulo $3$. Otherwise subtracting two choices that differ
in one sign would force all three coordinates to vanish. Modulo $3$ the
three rows of $L_0$ give $x-y-z$, its negative, and itself. Thus, after
choosing signs, applying $L_0$ makes all three coordinates nonzero modulo
$3$, preserves their oddness, and multiplies the squared norm by $9$.
Repeat $e$ times. If $e=0$, the sum of the three squares is $0$ modulo
$3$, and the fact that the coordinates are not all divisible by $3$
already makes all three nonzero modulo $3$. Finally choose signs so that
each coordinate is $-1$ modulo $6$. We have
\[
 M=(6u_1-1)^2+(6u_2-1)^2+(6u_3-1)^2,
\]
as required.
\end{proof}

This is the classical theorem that every nonnegative integer is a sum of
three generalized pentagonal numbers, using
$24\,x(3x-1)/2+1=(6x-1)^2$. Guy's observation and the norm-$9$
lifting identity are discussed and proved in Sun~\cite{sunpent}; see also
Oh~\cite{oh}. Neither that theorem nor the lifting identity is claimed as
new. Its uniformly nonzero residue is the feature exploited in the present
five-root merger.

\section{Soundness: every new zero restores a source zero}\label{sec:soundness}
Assume $P_{13}=0$ at an arbitrary signed integer assignment. The preceding
valuation argument supplies signed rational roots $t,f,d,t_+,y$ satisfying
\begin{gather}
 t^2=A_\tau,\quad f^2=D_f,\quad d^2=D_m,\quad t_+^2=S,
 \quad y^2=B_y,\label{eq:rootconditions}\\
 D_0+H(\xi+27V)=3Ht+9Hf+d+81Ht_++27Hy.\label{eq:carriereq}
\end{gather}
The signs in this equation are absorbed into the roots. Since their first
four squares are integers, $t,f,d,t_+$ are integers.

\subsection{Recover positivity and exclude every singular branch}
By Lemma~\ref{lem:pentagonal}, $N>0$; by Lemma~\ref{lem:joint}, $G>0$
and $Z=\eta\zeta>0$. In particular $k=\eta+\zeta$ is nonzero, its
absolute value is at least $2$, and $0<\eta/k<1$.

The literal affine forms exclude $q=0$ and $W=0$. Moreover $q$ is even,
$W$ is a nonzero multiple of $4$, and $Y=(2v+1)q^3$ is a nonzero
multiple of $8$. Hence $|X|\ge8$ and $|Y|\ge8$. From
$c/k=Y+\eta/k$ we obtain $|c|>2$. Also $a=(X+1)Y$ is a multiple of
$24$, so $\Delta=(a+2)^2-1\ge3$ is odd. Since $2r+1\ne0$,
\[
 Q=\Delta^2(2r+1)^2c^4>1.
\]
All these facts precede the rational-Pell integrality step.

\subsection{Rational auxiliary roots are integral}\label{sec:descent}
\begin{lemma}[Adjacent-coefficient Pell descent]\label{lem:descent}
If $Q>1$ and $V$ are integers and $y\in\Q$ satisfies
\[
 QV^2-(Q-1)y^2=1,
\]
then $y\in\Z$.
\end{lemma}
\begin{proof}
Change signs to take $V,y>0$. The equation excludes $V=0$. Write
$y=b/d$ in lowest terms with $d>0$. Multiplying by $d^2$ shows
$d^2\mid Q-1$. For $V=1$ the equation gives $y=1$.
If $V>1$, define
\[
 V'=(2Q-1)V-2(Q-1)y,\qquad y'=(2Q-1)y-2QV.
\]
Then $V'$ is an integer, the same norm equation holds, and the reduced
denominator of $y'$ is still $d$, because $2Q-1\equiv1\pmod d$.
The original equation gives $y>V$ and $QV>(Q-1)y$, whence
$-V<V'<V$. The norm excludes $V'=0$. Induction on $|V|$ reaches
$|V|=1$ while preserving the denominator, proving $d=1$.
\end{proof}
This argument is inherited from the rational Pell merger in package 69;
it is reproduced here because its exact hypotheses matter.
Apply it to~\eqref{eq:rootconditions}; thus $y$ is an integer as well.

\subsection{Explicit restoration maps}
Equation~\eqref{eq:carriereq} now implies $H\mid d-D_0$. Define
\begin{align}
 \gamma_{\rm old}&=\xi+27V-3t-9f-81t_+-27y,\notag\\
 I_{\rm old}&=\Delta I_0,
 &r_{\rm old}&=\frac{\Delta(2r+1)-1}{2},\label{eq:restore}\\
 T_{\rm old}&=fT,
 &y_{\rm gap}&=y-V,\notag\\
 \tau_{\rm half}&=2t-(2\ell+k),
 &J_{3,\rm old}&=2j.\notag
\end{align}
The only displayed division has an even numerator because $\Delta$ is
odd. Every restored coordinate is integral.

The first Pell residual is restored by the identity
\[
 \tau_{\rm half}^2+2(2\ell+k)\tau_{\rm half}+k^2-4
 =4\bigl(t^2-\ell(\ell+k)-1\bigr)=0.
\]
The main root is $d=X+ac+H\gamma_{\rm old}$ and satisfies
$d^2-\Delta c^2=1$. The normalized strong square gives
\[
 I_{\rm old}^2=\Delta(f^2-1),\qquad
 I_{\rm old}=(2r_{\rm old}+1)c^2.
\]
The CRT auxiliary expression becomes exactly
\[
 fcT_{\rm old}-f^2R-c=f^2(cT-R)-c=V.
\]
The auxiliary norm gives the parent's gap residual because
\[
 V^2+(1-I_{\rm old}^2)y_{\rm gap}(2V+y_{\rm gap})-1
 =I_{\rm old}^2V^2-(I_{\rm old}^2-1)y^2-1=0.
\]

Finally $N>0$ admits a representation
\[
 N=1+p_1^2+p_2^2+p_3(p_3+1).
\]
For example apply the three-square theorem to $4N-3\equiv1\pmod4$,
put its unique odd coordinate last, and divide the other coordinates by
$2$. This restores the exact parent's final positivity row.
All other parent definitions are literal polynomial identities. Thus every
new zero restores a zero of the seventeen-integer CRT parent, whose
source-to-refutation calibration is retained.

\section{Completeness: an explicit normalized canonical lift}\label{sec:complete}
Soundness was a restoration of arbitrary signed zeros. Completeness uses
the canonical witnesses supplied by the inherited native theorem. These
are different logical operations: the new restriction need not preserve
every old zero on the same coordinates.

\subsection{The inherited canonical interface}
A finite refutation supplies a source34 accepting tuple with
\[
 q=2^t,\quad t\text{ odd},\quad R\ge3,\quad R\equiv3\pmod4,
 \quad X=2^R,
\]
and $Y=(2v+1)q^3>0$, together with $G>0$, $Z>0$, the first Pell
root and the main root. Write $A_m=a+2$ and
\[
 d=\chi_{A_m}(R),\qquad c=\psi_{A_m}(R),
 \quad\chi_A(n)+\psi_A(n)\sqrt{A^2-1}=(A+\sqrt{A^2-1})^n.
\]
The determinant guard's support estimates are uniform over this same
canonical family. The main-root congruence gives an integral
$\gamma_{\rm old}$.

The canonical native repunit $J$ is odd and $1$ modulo $3$, hence $1$
modulo $6$. Thus $J=6j+1$ loses no canonical witness. The canonical
$W/4=2^{R-t-2}$ has even nonnegative exponent and is $1$ modulo $3$,
so $W=4(3w+1)$ remains valid. These facts are explicitly recorded in the
retained source70 orientation proof. Since $q$ is even, $a$ and $A_m$
are even, while $c$ is positive and odd.

\subsection{Construct the odd normalized strong witness}
Set
\[
 m=Rc,\quad I_0=\psi_{A_m}(m),\quad f=\chi_{A_m}(m).
\]
We claim $I_0/c^2$ is a positive odd integer. Pell composition gives
\[
 I_0=c\,\psi_d(c),\qquad d^2\equiv1\pmod c.
\]
The recurrence for $\psi$ proves
$\psi_d(n)\equiv n d^{n-1}\pmod c$, so $c\mid\psi_d(c)$.
Therefore $c^2\mid I_0$. Modulo $2$, the recurrence at even $A_m$
gives $\psi_{A_m}(n)\equiv n$. Both $m$ and $c$ are odd, so the
quotient $i=I_0/c^2$ is odd. Put $r=(i-1)/2$. The Pell identity now
gives $f^2=1+\Delta I_0^2=D_f$.

This is an explicit choice of a new strong witness. It does not assume that
the old one happened to have the normalization needed for compression.

\subsection{Construct the squared-modulus auxiliary witness}
Put $M=\Delta I_0$ and choose
\[
 V^{\dagger}=\frac{\chi_M(R)}{M},\qquad y=\psi_M(R).
\]
For odd $R=2b+1$, the quotient is an integer polynomial $Q_b(M^2)$.
The standard odd-index identities are
\[
 Q_b(0)=(-1)^bR,\qquad Q_b(1-A_m^2)=(-1)^b\psi_{A_m}(R).
\]
Since $R\equiv3\pmod4$, $b$ is odd. Also
$M^2=\Delta(f^2-1)\equiv1-A_m^2\pmod{f^2}$ and $c\mid M$.
Consequently
\[
 V^{\dagger}\equiv-c\pmod{f^2},\qquad
 V^{\dagger}\equiv-R\pmod c.
\]
Because $f^2\equiv1\pmod c$, the positive number
\[
 T=\frac{V^{\dagger}+c+Rf^2}{cf^2}
\]
is an integer. It makes $V=D_f(cT-R)-c=V^{\dagger}$. The Pell
identity at $M$ gives $QV^2-(Q-1)y^2=1$.
This squared-modulus argument is also supplied in source61, Section 5 of
its signed-index proof. It is independent of the native coefficient tuple.

\subsection{Choose the positivity and carrier coordinates}
The pair $G,Z>0$ permits an integer $s$ with $N>0$ by
Lemma~\ref{lem:joint}. Lemma~\ref{lem:pentagonal} then supplies
$u_1,u_2,t_+$ with $t_+^2=S$. Retain the canonical first Pell root $t$.
Set
\[
 \xi=\gamma_{\rm old}+3t+9f+27y+81t_+-27V.
\]
This is an integer and makes~\eqref{eq:carriereq} hold. One factor of
the five-root norm is zero, so $P_{13}=0$. Every new coordinate has been
chosen explicitly or by a proved integer representation theorem. This
completes both directions of Theorem~\ref{thm:main}.

\section{Exact degree and the fourteen-variable companion}\label{sec:degree}

\subsection{Uncentered degree}
Assign degree one to each independent integer coordinate and degree zero to
the six numerical ports. Direct substitution gives
\begin{center}\small
\begin{tabular}{@{}lrrrrrrrrrrrr@{}}\toprule
Register & $F$ & $R$ & $Y$ & $E,a,H$ & $k$ & $\ell$ & $I_0$ & $D_f$ & $Q$ & $V$ & $Z$ & $S$\\\midrule
Degree &2&5&4&6&7&17&3&18&30&23&22&44\\\bottomrule
\end{tabular}
\end{center}
In the rational norm, the five weighted roots have degrees at infinity
$23,15,7,28,29$ respectively. The last is uniquely largest. Thus the
uncentered norm has degree $32\cdot29+16\cdot30=1408$ after clearing
the auxiliary denominator. Equivalently, the highest term is $C_*^{16}$,
with $\deg C_*=88$.

\subsection{The degree-1392 shear}
Centering the auxiliary root is the integral coordinate change
$\xi\mapsto\xi+27V$. Since
\[
 B_y=V^2+\frac{V^2-1}{Q-1},
\]
one formal branch has $\sqrt{B_y}=V+O(t^{-7})$ after all variables
are scaled by $t$. Its weighted correction $27H(\sqrt{B_y}-V)$ has
degree at most $-1$. Of the 32 factors, 16 retain the leading auxiliary
term of degree 29, and 16 have it canceled. In the latter factors the
positivity root, of degree 28, is uniquely dominant. Therefore
\[
 \deg P_{13}\le16\cdot29+16\cdot28+16\cdot30=1392.
\]
This argument may be read in a formal Laurent-series extension of the
rational function field in the homogeneous input coordinates. It is a
degree calculation, not an approximation used in the soundness proof.

For an explicit nonvanishing leading homogeneous form, let $H_0,V_0,Q_0,S_0$
denote the leading homogeneous parts. Then the leading part of $P_{13}$ is
\begin{equation}\label{eq:leading}
 Q_0^{16}(54H_0V_0)^{16}(6561H_0^2S_0)^8.
\end{equation}
Each factor is nonzero. For example, if $k_0$ and $Y_0$ denote leading
parts, then $Z_0=-k_0^2Y_0^2$ and $S_0=-24Z_0^2\ne0$.
Similarly $V_0=-(D_f)_0R_0\ne0$. This proves equality in the degree
bound after the actual nonzero coefficient specialization. Independent
exact univariate polynomial calculations modulo a prime check the stated
degree and a nonzero leading coefficient of the delivered circuit.

\subsection{Fourteen integers and degree 652}
Keep three positivity coordinates $p_1,p_2,p_3$ in place of $u_1,u_2$,
and omit the fifth radicand. There are now fourteen unknowns. Define
\[
 G_0=N-\Pos(p_1,p_2,p_3),\qquad
 \Pos(p_1,p_2,p_3)=1+p_1^2+p_2^2+p_3(p_3+1).
\]
Use the carrier
\[
 p_{14}=2D_0+H\{\xi+6\ell+3k+18(a+2)I_0+54V\}.
\]
Let
\[
 N_{14}=\delta^8\mathcal F_4\left(
 p_{14};36H^2A_\tau,324H^2D_f,4D_m,
 \frac{2916H^2(QV^2-1)}\delta\right).
\]
The same Taylor prescription as~\eqref{eq:clear}, now with the three-root
norm of degree eight, specifies $N_{14}$ as an integer polynomial.
The companion is
\begin{equation}\label{eq:14}
 \boxed{\quad (1+4G_0)(1+4N_{14})-1=0.\quad}
\end{equation}
Both factors are integers congruent to $1$ modulo $4$. Their product can
equal $1$ only when both equal $1$. Thus this equation is exactly
$G_0=N_{14}=0$.

The four roots have distinct valuations, as before; multiplying all weights
by $2$ does not change that. The carrier recovers
$H\mid2(d-D_0)$, hence $H\mid d-D_0$ because $H$ is odd.
Conversely an old source tuple supplies an integer $\xi$ by the displayed
carrier identity. No new parity restriction on the free carrier is needed.

Centering the auxiliary, first, and normalized strong roots gives a norm of
exact degree 608. Its 16 branches split into groups of sizes $8,4,2,2$,
with respective degrees $29,23,15,7$, and denominator contribution $8\cdot30$:
\[
 8\cdot29+4\cdot23+2\cdot15+2\cdot7+8\cdot30=608.
\]
The guard residual has degree 44. Their unit-product equation therefore
has degree $608+44=652$. The delivered degree proof records its nonzero
leading homogeneous part, and the independent polynomial replay checks
the actual specialized circuit. The graph uses 230 operations.

The 748-, 700-, and 664-degree companions retain successively fewer
centering operations. They provide explicit control points for both the
algebra and the degree/cost tradeoff, rather than unsupported extrapolations.

\subsection{Fifteen integers and degree 214}
The archive also supplies a shorter three-root equation that keeps $f$ as
an independent coordinate and the ordinary strong row
$I^2-\Delta(f^2-1)=0$, where $I=(2r+1)c^2$.
Here $J=3j+1$, $Q=I^2$, and $V=fcT-f^2R-c$; all other native and
guard abbreviations agree with the parent. Set
\begin{gather*}
 p=2D_0+H\{\xi+3(2\ell+k)+18V\},\\
 A=36H^2A_\tau,\quad D=4D_m,\quad B_*=324H^2(QV^2-1),\quad\delta=Q-1,\\
 U=\delta(p^2-A-D)-B_*,\qquad
 L=U^2-4\delta^2AD-4\delta(A+D)B_*,\\
 N_{15}=L^2-64\delta^3ADB_*p^2.
\end{gather*}
The single equation is
\[
 (1+4\{I^2-\Delta(f^2-1)\})(1+4G_0)(1+4N_{15})-1=0.
\]
Its factors force the strong row, the guard and the three-root norm to
vanish. The guard supplies $Q>1$ before the norm is interpreted; the
strong row supplies $3\mid I$. Fixed valuations then recover the first,
auxiliary and main roots. Since $H$ is odd, the doubled carrier restores
the main congruence exactly. This is the independently reviewed initial
reduction, with no strong normalization or pentagonal substitution.
The norm, strong row and guard have degrees 156, 14 and 44, so the final
degree is 214. The complete derivation and literal 139-operation circuit
are retained in \file{forms/COMPACT15_PROOF.md} and its companion packet.

\section{Coefficients, verification, and logical scope}\label{sec:scope}

\subsection{The six fixed integers}
The archive's \file{coefficients/BINDING.json} is authoritative for names,
byte encoding and hashes. Integer bytes are minimal unsigned little-endian
representations; their compressed wrappers are not part of the mathematical
identity.
\begin{center}\small
\begin{tabular}{@{}lrrl@{}}\toprule
Port & Bits & Raw bytes & Definition\\\midrule
$\beta$ &1,064,921&133,116&$2^{1064921}-1$\\
$K_0$ &1,065,043&133,131&source34 native constructor\\
$M_C$ &1,064,921&133,116&source34 native constructor\\
$M_F$ &1,064,922&133,116&source34, already shifted\\
$A_0$ &1,114&140&$2^{1113}$\\
$B_0$ &666,105&83,264&$2^{666104}$\\\midrule
Total &&615,883&six fixed numeral payloads\\\bottomrule
\end{tabular}
\end{center}
The native recipe uses a literal 430-action binary machine table and its
943 local tiles. It constructs finite coefficient supports at radix
$\rho=2^{121}$, block length 8801, and outer radix
$B=\rho^{8801}=2^{1064921}$. In its notation,
\[
 K_0=D_C+B\rho,\qquad \beta=B-1,
 \qquad M_F=M_{F,\mathrm{native}}+\beta.
\]
The complete supports, padding, masks and action bytes are supplied, so the
recipe is finite and executable without simulating the machine. The final
verifier reconstructs all six integers and compares the raw values to their
payloads. Source31 and source34 share the same radix exponent but different
remaining ports; compact44 is a third tuple. No port substitution between
them is made here.

The new equations use the same six large numeral leaves as the calibrated
determinant parent, together with small integer multipliers. This does not
imply that the fully expanded coefficients have the same maximum height.
We do not claim a measured improvement in expanded monomial count or maximum
coefficient size. The polynomial is delivered as an exact finite arithmetic
graph, not as an enormous expanded list.

\subsection{What the executable checks establish}
The supplied standard-library verifier checks:
\begin{enumerate}
\item all actual fixed coefficient bytes and their deterministic reconstruction;
\item literal graph topology, permitted operations, unknown lists, live-node
      counts, and agreement with the finite construction recipes;
\item the generic denominator-cleared norm templates against a separately
      written exact multiquadratic-ring evaluator, including signed and
      rational radicands and explicit roots;
\item exact specialized polynomial degrees and nonzero leading coefficients,
      with the universal upper bounds justified in the written proof;
\item the normalized Pell, positivity and integral restoration identities,
      plus explicitly bounded diagnostic cases;
\item integrity of the packaged source files and negative controls showing
      that coefficient or graph corruption is detected.
\end{enumerate}
These checks do not search for a zero of the consistency polynomial and do
not test whether ZFC is consistent. Finite computations do not prove the
unbounded native encoding or the universal arithmetic lemmas. The latter
are supplied as written proofs and have separate internal mathematical
reviews.

\subsection{Proof status and dependency chain}
The logical composition is
\[
 \text{new integer zero}
 \longleftrightarrow\text{calibrated source34 zero}
 \longleftrightarrow\text{finite K215 accepting run}
 \longleftrightarrow\text{finite ZFC refutation}.
\]
The first equivalence is proved here, with explicit signed soundness and
canonical completeness. The later two are the named inherited ordinary
native/compiler and literal-machine calibration theorems. Their exact
source identities are recorded in \file{SOURCE_PINS.json}; the immediate
parent packet and its source archive are included.

All new arguments use finite field extensions, integer valuations, Pell
recurrences and classical integer representation theorems, and can be
formalized in ZFC. Thus the inherited ordinary calibration composes to the
ordinary theorem~\eqref{eq:target}. This report does not claim that the
optimized endpoint has been compiled in Lean or Rocq, or that an explicit
derivation has been exported in ProveIt's encoded object-language ZFC
calculus. The latter distinction is real and is preserved in the retained
formalization-status documents. The review here is AI-assisted internal
review, not external peer review.

\subsection{Reproduction}
From the extracted archive root run
\begin{verbatim}
python3 verify.py
python3 -O verify.py
\end{verbatim}
Both modes use the Python standard library. No network, machine-halting
search, numerical square-root approximation or proof assistant is required
for this finite replay. Output reports state their exact scope.
To rebuild the article, from \file{report/} run twice:
\begin{verbatim}
pdflatex -interaction=nonstopmode -halt-on-error report.tex
\end{verbatim}
The release PDF was compiled and visually inspected separately from the
arithmetic replay. The ZIP contains complete LaTeX sources, the PDF, literal
equations, numerical coefficients, construction and verification code, proof
notes, independent reviews and provenance.

\section{What this changes, and what remains open}
The practical contribution is a new modest-degree range for equations over
\emph{unrestricted integers}. The inspected small-degree chain previously
ended at seventeen variables and degree 88; the new points at fourteen and
thirteen variables have degrees 652 and 1392. The huge-degree thirteen-
integer construction is no longer the only delivered thirteen-integer
option in this research line.

The reusable idea is a local arithmetic invariant: make every relevant
radicand a unit at one prime, then assign distinct valuations through fixed
weights. It removes the need to separate every possible signed subset by
large real inequalities. The strong normalization and positivity conversion
are useful precisely because they preserve this invariant. Applying the
combiner without those uniform residue checks would be unsound.

The contribution does not lower the ten-integer or nine-positive general
variable counts, solve the consistency question, establish worldwide
priority, or prove that any delivered equation is globally simplest. It
also retains substantial inherited coefficient data. A further improvement
should measure variables, degree, live circuit size and coefficient size
separately, and should preserve the explicit endpoint calibration.

The next concrete formalization target is now small and well delimited:
the valuation-separated square-class lemma, the normalized odd-index Pell
adapter, the $3$-adic-unit positivity lemma, and the exact five-root norm
identity. None requires redoing the native machine compiler. Completing
those kernel proofs would strengthen verification of the new reduction;
the native source-to-refutation internalization remains its own task.

\appendix
\section{The immediate seventeen-integer parent}
The literal parent is retained as
\file{upstream/transport/degree88_17.json}. It uses the same native and
guard abbreviations, with $J=3J_3+1$, ordinary strong variable $f$, and
three separate root coordinates. Put
\[
 I=(2r+1)c^2,\quad V=fcT-f^2R-c,
 \quad D=X+ac+H\gamma,\quad P=\ell+k.
\]
Its five residuals are
\begin{align*}
 &\tau^2+2\tau(2P-k)+k^2-4,\\
 &D^2-\Delta c^2-1,\\
 &I^2-\Delta(f^2-1),\\
 &V^2+(1-I^2)y(2V+y)-1,\\
 &GZ-(s^2-G-Z+1)^2-\Pos(p_1,p_2,p_3).
\end{align*}
The parent equation is their sum of squares. Its unknowns are
$J_3,C,z,f,h,r,v,w,\tau,y,\gamma,T,c,s,p_1,p_2,p_3$.
The definitions in Sections~\ref{sec:soundness}--\ref{sec:complete}
give the exact maps needed to compare this parent with the new equations.

\section{The source boundary is a mathematical interface}
The retained parent uses the calibrated determinant condition $G>0$ and
the interval condition $\eta\zeta>0$. Soundness of the determinant applies
to arbitrary signed assignments; completeness uses fixed leading-digit
bounds of canonical native witnesses. The residue orientation imposes
$X\equiv2\pmod3$ to select the positive signed power of two. The CRT
compression of auxiliary quotients uses the strong rank lemma before the
auxiliary rank is identified, avoiding a circular use of the target rank.
The original proofs of these facts are retained with the parent archive.

The present reduction does not strengthen any of these inherited
interfaces silently. The only additional canonical restrictions are
$J\equiv1\pmod6$ and the normalized odd strong witness. The first follows
from the already stated canonical repunit; the second is constructed in
Section~\ref{sec:complete}. The pentagonal representation holds for every
positive guard value. All compatibility claims therefore have explicit
proofs.

\begin{thebibliography}{9}
\bibitem{repo} V. Reshetnikov and contributors, \emph{ProveIt}, comparison
commit\\\pin, the DiophantineConsistency collection, active formalization
documents and incoming archives.
\url{https://github.com/VladimirReshetnikov/ProveIt/tree/cfbf118b4ad9b2c651fe96b571d9415261766328/Computability/DiophantineConsistency}.

\bibitem{parent} \emph{Pell congruences and Diophantine consistency},
incoming archive \file{con_zfc_pell_congruences_2026-10-23.zip}, including
the degree-88 seventeen-integer CRT/transport parent. Original archive,
proofs, coefficient constructor and source notices retained in this release.

\bibitem{native} ProveIt packages 34, 61, 69, 70 and 71:
source34 mixed-parity native coefficients; signed auxiliary-index and
squared-modulus reconstruction; rational Pell descent; residue orientation;
calibrated determinant guard. All referenced source identities are pinned
in the release provenance.

\bibitem{sunpent} Z.-W. Sun, \emph{On universal sums of polygonal numbers},
Science China Mathematics \textbf{58} (2015), 1367--1396.
\url{https://arxiv.org/pdf/0905.0635}. In particular the discussion of Guy's
three-generalized-pentagonal observation and the norm-$9$ identity (1.4).

\bibitem{oh} B.-K. Oh, \emph{Ternary universal sums of generalized
pentagonal numbers}, Journal of the Korean Mathematical Society
\textbf{48} (2011), 837--847. \url{https://arxiv.org/pdf/0911.1181}.

\bibitem{mr} Y. Matiyasevich and J. Robinson, \emph{Reduction of an arbitrary
Diophantine equation to one in 13 unknowns}, Acta Arithmetica \textbf{27}
(1975), 521--553. \url{https://doi.org/10.4064/aa-27-1-521-553}.
Classical background for relation combining; the fixed-prime specialization
used here is proved directly in Section~\ref{sec:valuation}.
\end{thebibliography}
\end{document}
